\documentclass[11pt,a4paper]{article}

% ============================================================
% PACKAGES
% ============================================================
\usepackage[utf8]{inputenc}
\usepackage[T1]{fontenc}
\usepackage[english]{babel}
\usepackage{amsmath,amssymb,amsthm,mathtools}
\usepackage{mathrsfs}
\usepackage{tikz-cd}
\usepackage{enumitem}
\usepackage{hyperref}
\usepackage{geometry}
\usepackage{microtype}
\usepackage{booktabs}
\usepackage{array}
\usepackage{longtable}
\usepackage{graphicx}

\geometry{margin=1in}

% ============================================================
% THEOREM ENVIRONMENTS
% ============================================================
\theoremstyle{definition}
\newtheorem{definition}{Definition}[section]
\newtheorem{axiom}[definition]{Axiom}
\newtheorem{example}[definition]{Example}
\newtheorem{remark}[definition]{Remark}
\newtheorem{convention}[definition]{Convention}
\newtheorem{construction}[definition]{Construction}
\newtheorem{assumption}[definition]{Assumption}
\newtheorem{protocol}[definition]{Protocol}

\theoremstyle{plain}
\newtheorem{proposition}[definition]{Proposition}
\newtheorem{theorem}[definition]{Theorem}
\newtheorem{lemma}[definition]{Lemma}
\newtheorem{corollary}[definition]{Corollary}

% ============================================================
% MACROS (inherited from framework paper)
% ============================================================
\newcommand{\Lam}{\Lambda}
\newcommand{\Sp}{\mathrm{Sp}}
\newcommand{\Ge}{\mathrm{Ge}}
\newcommand{\To}{\mathrm{To}}
\newcommand{\Al}{\mathrm{Al}}
\newcommand{\Nu}{\mathrm{Nu}}
\newcommand{\Fo}{\mathrm{Fo}}
\newcommand{\Cl}{\mathrm{Cl}}
\newcommand{\Corr}{\mathrm{Corr}}
\newcommand{\Inv}{\mathrm{Inv}}
\newcommand{\An}{\mathrm{An}}
\newcommand{\Ob}{\mathfrak{O}}
\newcommand{\Ecal}{\mathcal{E}}
\newcommand{\Ocal}{\mathcal{O}}
\newcommand{\Acal}{\mathcal{A}}
\newcommand{\Ccal}{\mathcal{C}}
\newcommand{\Mcal}{\mathcal{M}}
\newcommand{\Ncal}{\mathcal{N}}
\newcommand{\Rcal}{\mathcal{R}}
\newcommand{\Dcal}{\mathcal{D}}
\newcommand{\philam}{\phi_{\lambda\mu}}
\newcommand{\Vmat}{\mathbf{V}}
\newcommand{\Jset}{\mathcal{J}}
\newcommand{\dd}{\,\mathrm{d}}
\newcommand{\op}[1]{\operatorname{#1}}
\newcommand{\spec}{\sigma}
\newcommand{\Norm}[1]{\left\lVert #1 \right\rVert}
\newcommand{\Abs}[1]{\left\lvert #1 \right\rvert}
\newcommand{\ip}[2]{\left\langle #1, #2 \right\rangle}
\newcommand{\R}{\mathbb{R}}
\newcommand{\C}{\mathbb{C}}
\newcommand{\N}{\mathbb{N}}
\newcommand{\Z}{\mathbb{Z}}
\newcommand{\Pmark}{\mathrm{P}}
\newcommand{\Nmark}{\mathrm{N}}
\newcommand{\Omark}{\mathrm{O}}
\newcommand{\Fmark}{\mathrm{F}}
\newcommand{\Fr}{\mathrm{Fr}}

% ============================================================
% TITLE
% ============================================================
\title{%
  \textbf{Structural Constraining and Spectral Emergence:}\\[4pt]
  \large A Multirepresentational Methodology for Nonlinear Evolution Equations
}
\author{}
\date{}

\begin{document}
\maketitle

\begin{abstract}
\noindent
We develop a methodology for the analysis of mathematical models by
successive representation across spectral, geometric, topological,
algebraic, numerical, and formal layers. Unlike standard
multi-physics or multi-scale approaches, the methodology treats each
layer as a \emph{constraint} on the space of admissible models,
rather than as an independent representation. We introduce three
structural principles: \emph{constraining}, according to which the
solution space shrinks monotonically under layer transitions;
\emph{spectral emergence}, according to which the spectral layer
may be constructed rather than assumed, via Fr\'echet linearisation
of a nonlinear operator at an equilibrium; and
\emph{discrepancy classification}, according to which any mismatch
between the initial model and its reconstruction falls into one of
three regimes---discretisation, physical extension, or structural
inconsistency. Four worked examples, all elementary enough to be
verified by hand, exhibit the full methodology: the Dirichlet
Laplacian on $[0,1]$, the path-vs-cycle bifurcation, an erroneous
diffusion equation with spurious advection, and a nonlinear
$p$-Laplacian gradient flow. For each example we produce a
validation matrix recording the epistemic status of each layer. We
close by discussing the limits of the methodology, its relationship
to structure-preserving numerics, inverse problems, and
proof-assistant formalisation, and open problems.
\end{abstract}

\tableofcontents

% ============================================================
\part{Foundations}
% ============================================================

\section{Introduction}
\label{sec:intro}

\subsection{Motivation}

Given a mathematical model---typically a partial differential
equation, an operator equation, or a variational principle---the
question of \emph{what structure the model carries} is rarely
answered by the model's surface form. A diffusion equation hides a
gradient-flow structure; a linear equation hides a spectral
decomposition valid only near an equilibrium; an apparent scalar
equation hides a symmetry group; and so on.

The methodology developed here is a systematic procedure for
exposing such hidden structure. Its central claim is that
mathematical representations should be treated not as
\emph{parallel descriptions} of a fixed object, but as
\emph{constraints} that progressively eliminate candidate models.

\subsection{Relationship to the framework paper}

This paper inherits notation and definitions from a companion
framework paper \cite{framework} which introduces the master object
$\Ob$, the layer set $\Lam$, the transition morphisms $\philam$, the
reconstruction operators $\rho_\Cl, \rho_\Inv, \rho_\An$, and the
validation matrix $\Vmat$. The reader is referred to \cite{framework}
for the full treatment of the six construction layers and the three
reconstruction layers. The present paper assumes this vocabulary and
develops three new structural principles.

\subsection{Contributions}

\begin{enumerate}[label=(\roman*)]
  \item \textbf{Structural constraining} (\S\ref{sec:constraining}):
        a formal definition of the methodology as a monotone
        intersection of admissible model spaces.
  \item \textbf{Spectral emergence} (\S\ref{sec:emergence}): a
        construction by which the spectral layer is produced by
        Fr\'echet linearisation of a nonlinear operator, rather than
        assumed.
  \item \textbf{Discrepancy classification} (\S\ref{sec:discrepancy}):
        a trichotomy of reconstruction discrepancies into
        discretisation, physical extension, and structural
        inconsistency.
  \item \textbf{Four worked examples} (\S\ref{sec:toys}) covering
        linear, bifurcating, erroneous, and nonlinear cases.
  \item \textbf{Validation matrices} (\S\ref{sec:validation-new})
        for each example, recording the epistemic status of each
        layer.
\end{enumerate}

\subsection{Organization}

Part~I (\S\ref{sec:constraining}--\ref{sec:discrepancy}) develops the
three structural principles. Part~II (\S\ref{sec:toys}) presents the
four worked examples. Part~III (\S\ref{sec:validation-new}--\ref{sec:discussion})
treats the validation matrix, the comparison with related work, and
the limits of the methodology.

% ------------------------------------------------------------
\section{Structural Constraining}
\label{sec:constraining}
% ------------------------------------------------------------

\subsection{The Constraining Space}

Let $\Ob$ be the master object of the framework paper and let
\[
  \Lam = \{\Sp, \Ge, \To, \Al, \Nu, \Fo, \Cl, \Corr, \An\}
\]
be the extended layer set. Fix a problem class $\mathcal P$: the set
of models (operators, equations, variational principles) compatible
with the empirical and formal data under consideration.

\begin{definition}[Constrained model space]\label{def:constrained-space}
For each layer $\lambda \in \Lam$, the \emph{constrained model space
at layer $\lambda$} is
\[
  \Ccal_\lambda
  \;=\;
  \bigl\{ P \in \Ccal_{\lambda^-} \;:\;
    P \text{ is compatible with } \Ocal_\lambda \bigr\},
\]
where $\lambda^-$ is the immediate predecessor of $\lambda$ in the
construction order and $\Ccal_{\varnothing} = \mathcal P$.
\end{definition}

\begin{axiom}[Monotonicity of constraining]\label{ax:monotone}
For every pair $\lambda \prec \mu$ in $\Lam$,
\[
  \Ccal_\mu \;\subseteq\; \Ccal_\lambda .
\]
\end{axiom}

\begin{proposition}[Chain of constraints]\label{prop:chain}
The constrained model spaces form a decreasing chain
\[
  \mathcal P
  \;\supseteq\;
  \Ccal_\Sp
  \;\supseteq\;
  \Ccal_\Ge
  \;\supseteq\;
  \Ccal_\To
  \;\supseteq\;
  \Ccal_\Al
  \;\supseteq\;
  \Ccal_\Nu
  \;\supseteq\;
  \Ccal_\An .
\]
\end{proposition}

\begin{proof}
Immediate from Axiom~\ref{ax:monotone} and transitivity of $\prec$.
\end{proof}

\subsection{The Solution as an Intersection}

\begin{definition}[Constrained solution set]\label{def:sol-set}
The \emph{constrained solution set} of the methodology is
\[
  \Ccal_\ast
  \;=\;
  \bigcap_{\lambda \in \Lam} \Ccal_\lambda .
\]
\end{definition}

\begin{proposition}[Model determination]\label{prop:determination}
If $\Ccal_\ast$ is a singleton $\{P_\ast\}$, then the methodology
\emph{determines} the model $P_\ast$ uniquely.
\end{proposition}

\begin{proof}
By definition of the intersection.
\end{proof}

\begin{remark}
The methodology does \emph{not} require $\Ccal_\ast$ to be a
singleton. In many applications, $\Ccal_\ast$ has positive
dimension, corresponding to a moduli space of admissible models.
The methodology's output is then a parametrised family, and the
parameters carry physical or mathematical meaning (symmetries,
boundary conditions, coupling constants).
\end{remark}

\subsection{The Constraining Diagram}

The constraining principle can be displayed as
\begin{equation}\label{eq:constraining-diagram}
\begin{tikzcd}[row sep=large, column sep=large]
  \mathcal P \arrow[d, hookrightarrow] \\
  \Ccal_\Sp \arrow[d, hookrightarrow] \\
  \Ccal_\Ge \arrow[d, hookrightarrow] \\
  \Ccal_\To \arrow[d, hookrightarrow] \\
  \Ccal_\Al \arrow[d, hookrightarrow] \\
  \Ccal_\Nu \arrow[d, hookrightarrow] \\
  \Ccal_\An
\end{tikzcd}
\end{equation}
where each inclusion is the constraining imposed by the corresponding
layer.

\subsection{Comparison with the Framework Paper}

In the framework paper, the layers are treated as \emph{projections
of a single master object}; here, they are treated as
\emph{constraints on a model class}. The two perspectives are
complementary:
\begin{itemize}
  \item the projection view describes \emph{what each layer sees};
  \item the constraining view describes \emph{what each layer
        eliminates}.
\end{itemize}

\begin{proposition}[Equivalence of perspectives]\label{prop:equivalence}
Let $\pi_\lambda : \Ob \to \Ocal_\lambda$ be the projections of the
framework paper. Then
\[
  \Ccal_\lambda
  =
  \bigl\{ P \in \mathcal P : \pi_\lambda(P) \in \Ocal_\lambda \bigr\},
\]
where $\pi_\lambda(P)$ denotes the layer-$\lambda$ shadow of the
model $P$.
\end{proposition}

\begin{proof}
By definition of the projection and of the constrained model space.
\end{proof}

% ------------------------------------------------------------
\section{Spectral Emergence}
\label{sec:emergence}
% ------------------------------------------------------------

\subsection{The Problem with Linear Spectral Theory}

In the framework paper, the spectral layer $\Ocal_\Sp$ is defined
by the eigenpair equation
\[
  T v_\lambda = \lambda v_\lambda
\]
for a linear operator $T$. This presumes $T$ is linear. For
nonlinear evolution equations
\[
  u_t = \Ncal(u),
\]
no such decomposition exists globally.

\begin{definition}[Nonlinear evolution equation]\label{def:nonlinear-eq}
A \emph{nonlinear evolution equation} on a domain $\Omega$ is a
relation
\[
  u_t = \Ncal(u),
  \qquad u(\cdot, t) \in X,
\]
where $X$ is a function space and $\Ncal : D(\Ncal) \subset X \to X$
is a nonlinear operator.
\end{definition}

\begin{remark}
If $\Ncal$ is nonlinear, the set of eigenvalues and eigenfunctions
of $\Ncal$ is not well-defined. The spectral layer of the framework
paper is therefore not directly available.
\end{remark}

\subsection{Fr\'echet Linearisation}

\begin{definition}[Fr\'echet derivative]\label{def:frechet}
Let $u_\ast \in X$ be an equilibrium of \eqref{eq:nonlinear-flow},
i.e.\ $\Ncal(u_\ast) = 0$. The \emph{Fr\'echet derivative} of
$\Ncal$ at $u_\ast$ is the bounded linear operator
$D\Ncal[u_\ast] : X \to X$ such that
\[
  \Ncal(u_\ast + v) - \Ncal(u_\ast)
  =
  D\Ncal[u_\ast]\, v + o(\Norm{v}_X)
  \qquad (\Norm{v}_X \to 0).
\]
\end{definition}

\begin{construction}[Emergent spectral layer]\label{constr:emergent-spectral}
Given a nonlinear evolution equation $u_t = \Ncal(u)$ with
equilibrium $u_\ast$, the \emph{emergent spectral layer} is
\[
  \Ocal_\Sp^{\mathrm{emg}}
  \;=\;
  \bigl(
    L_{u_\ast},
    \spec(L_{u_\ast}),
    E_{L_{u_\ast}}
  \bigr),
\]
where $L_{u_\ast} := D\Ncal[u_\ast]$ is the Fr\'echet derivative of
$\Ncal$ at $u_\ast$.
\end{construction}

\begin{proposition}[Linearisation as first-order approximation]\label{prop:linearisation}
For $\Norm{v}_X$ small,
\[
  u_t = \Ncal(u_\ast + v)
  =
  L_{u_\ast} v + o(\Norm{v}_X).
\]
Therefore, the spectral data of $L_{u_\ast}$ govern the linear
dynamics of small perturbations around the equilibrium $u_\ast$.
\end{proposition}

\begin{proof}
Direct application of Definition~\ref{def:frechet}.
\end{proof}

\subsection{The Emergence Diagram}

The spectral layer is now constructed, not assumed:
\begin{equation}\label{eq:emergence-diagram}
\begin{tikzcd}[row sep=large, column sep=large]
  \Ncal \arrow[r, "\text{linearise at } u_\ast"] \arrow[dr, dashed]
    & L_{u_\ast} \arrow[d, "\text{spectral theorem}"] \\
  & \Ocal_\Sp^{\mathrm{emg}}
\end{tikzcd}
\end{equation}

\subsection{Limits of Emergent Spectral Data}

\begin{proposition}[Insufficiency of emergent spectrum]\label{prop:insufficient}
The emergent spectral layer $\Ocal_\Sp^{\mathrm{emg}}$ captures only
the \emph{linear} behaviour of $\Ncal$ near $u_\ast$. Higher-order
terms $\Ncal_k$, $k \geq 2$, in the Taylor expansion
\[
  \Ncal(u_\ast + v)
  =
  L_{u_\ast} v + \Ncal_2(v, v) + \Ncal_3(v, v, v) + \cdots
\]
are not recoverable from $\Ocal_\Sp^{\mathrm{emg}}$ alone.
\end{proposition}

\begin{proof}
The spectral theorem applies to the linear operator $L_{u_\ast}$
only. The higher-order terms are nonlinear and therefore lie outside
the domain of linear spectral theory.
\end{proof}

\begin{remark}
This limitation is precisely why the methodology requires
\emph{all} layers: the spectral layer constrains the linear part,
but only the algebraic, numerical, and formal layers recover the
full nonlinear structure.
\end{remark}

% ------------------------------------------------------------
\section{Discrepancy Classification}
\label{sec:discrepancy}
% ------------------------------------------------------------

\subsection{Reconstruction Discrepancies}

Given a candidate model $P_{\rm cand}$ and its reconstruction
$\Rcal(\Ecal(P_{\rm cand}))$ (notation inherited from the framework
paper), the \emph{discrepancy} is
\[
  \Delta P
  \;=\;
  \Rcal\bigl(\Ecal(P_{\rm cand})\bigr) - P_{\rm cand}.
\]

\begin{theorem}[Trichotomy of discrepancies]\label{thm:trichotomy}
Every discrepancy $\Delta P$ arising from the methodology falls into
exactly one of the following three regimes:
\begin{enumerate}[label=\textbf{(\roman*)}]
  \item \textbf{Discretisation regime:}
        $\Delta P = \varepsilon_N$ with $\varepsilon_N \to 0$ as
        $N \to \infty$. The discrepancy is a numerical artefact.
  \item \textbf{Extension regime:}
        $\Delta P = \Delta_\ast \neq 0$, independent of $N$. The
        candidate model was incomplete.
  \item \textbf{Inconsistency regime:}
        $\Delta P$ violates a structural constraint (symmetry,
        self-adjointness, topological invariant, energy law) required
        by the problem. The candidate model was inconsistent.
\end{enumerate}
\end{theorem}

\begin{proof}
Suppose $\Delta P \neq 0$. Consider the limit
$\Delta_\ast := \lim_{N \to \infty} \Delta P$.
\begin{itemize}
  \item If $\Delta_\ast = 0$, then $\Delta P \to 0$, giving (i).
  \item If $\Delta_\ast \neq 0$ and $\Delta_\ast$ is compatible with
        all layer constraints, then (ii) holds.
  \item If $\Delta_\ast \neq 0$ and $\Delta_\ast$ violates some
        layer constraint, then (iii) holds.
\end{itemize}
These cases are exhaustive and mutually exclusive.
\end{proof}

\subsection{Diagnostic Procedure}

\begin{protocol}[Discrepancy diagnosis]\label{prot:diagnosis}
Given a candidate model $P_{\rm cand}$ and its reconstruction
$\Rcal(\Ecal(P_{\rm cand}))$:
\begin{enumerate}[label=(\arabic*)]
  \item Compute $\Delta P$.
  \item If $\Norm{\Delta P} = \mathcal{O}(N^{-p})$ for some $p > 0$,
        classify as \emph{discretisation}. Terminate.
  \item Otherwise, check whether $\Delta P$ violates any of:
        \begin{itemize}
          \item self-adjointness: $L \neq L^\ast$;
          \item symmetry: $[L, S] \neq 0$ for a problem symmetry $S$;
          \item energy law: $\frac{dE}{dt} > 0$ for a dissipative
                equation;
          \item topological invariant: $H_k \neq H_k^{\rm target}$.
        \end{itemize}
        If any violation is found, classify as
        \emph{inconsistency}. Propose a correction
        $P_{\rm corr} = P_{\rm cand} - \Delta_{\rm structural}$.
  \item Otherwise, classify as \emph{extension}. Augment the model
        with $\Delta P$.
\end{enumerate}
\end{protocol}

\subsection{The Three Regimes Illustrated}

\begin{center}
\begin{tabular}{c|c|c}
  \toprule
  Regime & Signature & Interpretation \\
  \midrule
  Discretisation & $\Delta P \to 0$ as $N \to \infty$
    & Numerical artefact \\
  Extension & $\Delta P \to \Delta_\ast \neq 0$, compatible
    & Missing physics \\
  Inconsistency & $\Delta P$ violates structural constraint
    & Candidate was wrong \\
  \bottomrule
\end{tabular}
\end{center}

% ============================================================
\part{Worked Examples}
% ============================================================

\section{Four Toys}
\label{sec:toys}

Each toy below is elementary enough to be verified by hand, and each
exhibits the full methodology. Section~\ref{sec:toy1} treats the
Dirichlet Laplacian as the baseline case (spectrum available as
input). Section~\ref{sec:toy2} exhibits the geometry bifurcation
(path vs.\ cycle). Section~\ref{sec:toy3} exhibits error detection
via symmetry violation. Section~\ref{sec:toy4} exhibits spectral
emergence via Fr\'echet linearisation.

% ------------------------------------------------------------
\section{Toy 1: Dirichlet Laplacian on $[0,1]$}
\label{sec:toy1}
% ------------------------------------------------------------

\subsection{Candidate model}

\[
  P_{\rm cand}: \qquad -u''(x) = f(x), \quad x \in (0,1),
  \quad u(0) = u(1) = 0.
\]

\subsection{Spectral layer}

The operator $A = -\partial_{xx}$ on
$D(A) = H^2(0,1) \cap H_0^1(0,1)$ is self-adjoint, positive, and
unbounded. Its spectrum is
\[
  \spec(A) = \{n^2 \pi^2 : n \in \N, n \geq 1\}
\]
with eigenfunctions $\phi_n(x) = \sin(n \pi x)$.

\subsection{Geometric layer}

$M = [0,1]$ with metric $g = dx^2$, and $A = -\Delta_g$ on $M$.

\subsection{Topological layer}

$H_0(M) \cong \Z$, $H_1(M) = 0$, $\pi_1(M) = 0$.

\subsection{Algebraic layer}

$A = D^\ast D$ where $D = \partial_x$ on $H_0^1(0,1)$; therefore
$A \succeq 0$.

\subsection{Numerical layer}

Tridiagonal $A_N$ with $h = 1/(N+1)$; eigenvalues
\[
  \lambda_k^{(N)}
  =
  \frac{4}{h^2} \sin^2\!\left( \frac{k \pi}{2(N+1)} \right),
\]
satisfying $\lambda_k^{(N)} = k^2 \pi^2 + \mathcal{O}(h^2)$.

\subsection{Formal and reconstruction layers}

Formal theory $\Gamma$ encodes the weak formulation. The
reconstruction $\Rcal(\Ecal(P_{\rm cand}))$ recovers $P_{\rm cand}$.

\subsection{Discrepancy classification}

$\Delta P = \mathcal{O}(h^2) \to 0$. Regime (i): discretisation.

\subsection{Validation matrix}

\begin{center}
\begin{tabular}{l|cccccccc}
  \toprule
  Layer & $\mathrm{Ex}$ & $\mathrm{Co}$ & $\mathrm{In}$
  & $\mathrm{Nu}$ & $\mathrm{St}$ & $\mathrm{Cv}$
  & $\mathrm{Re}$ & $\mathrm{Rev}$ \\
  \midrule
  $\Sp$   & $\Pmark$ & $\Pmark$ & $\Pmark$ & $\Nmark$ & $\Pmark$ & $\Pmark$ & $\Pmark$ & $\Pmark$ \\
  $\Ge$   & $\Pmark$ & $\Pmark$ & $\Pmark$ & $\Nmark$ & $\Pmark$ & $\Pmark$ & $\Pmark$ & $\Pmark$ \\
  $\To$   & $\Pmark$ & $\Pmark$ & $\Pmark$ & $\Omark$ & $\Pmark$ & $\Pmark$ & $\Pmark$ & $\Pmark$ \\
  $\Al$   & $\Pmark$ & $\Pmark$ & $\Pmark$ & $\Omark$ & $\Pmark$ & $\Pmark$ & $\Pmark$ & $\Pmark$ \\
  $\Nu$   & $\Pmark$ & $\Pmark$ & $\Pmark$ & $\Nmark$ & $\Pmark$ & $\Pmark$ & $\Omark$ & $\Omark$ \\
  $\Fo$   & $\Pmark$ & $\Pmark$ & $\Pmark$ & $\Omark$ & $\Pmark$ & $\Pmark$ & $\Pmark$ & $\Pmark$ \\
  \bottomrule
\end{tabular}
\end{center}

% ------------------------------------------------------------
\section{Toy 2: Geometry Bifurcation (Path vs.\ Cycle)}
\label{sec:toy2}
% ------------------------------------------------------------

\subsection{The same local principle, two geometries}

Consider $N$ nodes with nearest-neighbour interaction. Two
geometric realisations are possible:
\begin{itemize}
  \item \textbf{Path:} $G_N^{(A)} = P_N$, approximating $M_A = [0,1]$.
  \item \textbf{Cycle:} $G_N^{(B)} = C_N$, approximating $M_B = S^1$.
\end{itemize}

\subsection{Constraining diagram}

\begin{equation}
\begin{tikzcd}[row sep=large, column sep=large]
  & \text{local interaction} \arrow[dl] \arrow[dr] & \\
  M_A = [0,1] \arrow[d] & & M_B = S^1 \arrow[d] \\
  \partial M_A \neq \varnothing \arrow[d] & & \partial M_B = \varnothing \arrow[d] \\
  \text{Dirichlet BC} \arrow[d] & & \text{periodic BC} \arrow[d] \\
  \sin(n \pi x) \arrow[d] & & e^{i k \theta} \arrow[d] \\
  \lambda_n = n^2 \pi^2 & & \lambda_k = k^2
\end{tikzcd}
\end{equation}

\subsection{Spectral comparison}

\begin{center}
\begin{tabular}{l|c|c}
  \toprule
  Feature & Path ($[0,1]$) & Cycle ($S^1$) \\
  \midrule
  $\pi_1$ & $0$ & $\Z$ \\
  Boundary & $\{0,1\}$ & $\varnothing$ \\
  Modes & $\sin(n \pi x)$ & $e^{ik\theta}$ \\
  Spectrum & $n^2 \pi^2$ & $k^2$ \\
  Multiplicity & simple & double for $k \neq 0$ \\
  Natural numerical method & finite differences & FFT \\
  \bottomrule
\end{tabular}
\end{center}

\subsection{Constraining interpretation}

Both realisations share the discrete object $L_N$. The geometric
layer selects one of two families; the topological layer then
restricts the admissible functional analysis (Dirichlet vs.\ periodic
Hilbert spaces); the algebraic layer determines the mode basis
(sine vs.\ Fourier); and the numerical layer is dictated by the
geometry. The two branches converge to genuinely different models
$P_{\rm cand}^{(A)} \neq P_{\rm cand}^{(B)}$.

\begin{proposition}[Non-equivalence of branches]\label{prop:non-equiv}
$P_{\rm cand}^{(A)} \not\cong P_{\rm cand}^{(B)}$ as operator
equations.
\end{proposition}

\begin{proof}
The spectra $\{n^2 \pi^2\}_{n \geq 1}$ and $\{k^2\}_{k \geq 0}$ are
not isomorphic as subsets of $\R$: the latter contains $0$ (constant
mode) and all non-zero eigenvalues have multiplicity $2$, whereas
the former has neither property.
\end{proof}

% ------------------------------------------------------------
\section{Toy 3: Erroneous Diffusion with Spurious Advection}
\label{sec:toy3}
% ------------------------------------------------------------

\subsection{Candidate model}

\[
  P_{\rm cand}: \qquad u_t = \kappa u_{xx} + \varepsilon u_x,
  \quad x \in (0,1), \quad u(0,t) = u(1,t) = 0.
\]

The candidate is \emph{assumed} to model isotropic diffusion, but
contains an additional advection term.

\subsection{Layer-by-layer analysis}

\begin{enumerate}[label=(\arabic*)]
  \item \textbf{Spectral:} removing the first-order term via
        $\phi(x) = e^{-\varepsilon x/(2\kappa)} v(x)$ gives
        \[
          \lambda_n = -\kappa n^2 \pi^2 - \frac{\varepsilon^2}{4\kappa},
        \]
        which departs from the pure diffusion spectrum by the
        constant shift $-\varepsilon^2/(4\kappa)$.
  \item \textbf{Geometric:} the term $\varepsilon \partial_x$
        selects a preferred direction on $M = [0,1]$; the geometry
        $(M, g)$ alone does not provide such a direction.
  \item \textbf{Algebraic:} decompose $L = L_S + L_A$ where
        $L_S = \kappa \partial_{xx}$ is symmetric and
        $L_A = \varepsilon \partial_x$ is antisymmetric. Then
        $L \neq L^\ast$.
  \item \textbf{Energy:} for pure diffusion,
        $\frac{d}{dt} \Norm{u}_{L^2}^2 = -2\kappa \Norm{u_x}_{L^2}^2$.
        For the candidate,
        \[
          \frac{1}{2}\frac{d}{dt} \Norm{u}_{L^2}^2
          =
          -\kappa \Norm{u_x}_{L^2}^2
          +
          \frac{\varepsilon}{2} [u^2]_0^1
          =
          -\kappa \Norm{u_x}_{L^2}^2
        \]
        because $u(0,t) = u(1,t) = 0$. \emph{The energy test does not
        detect the error.}
  \item \textbf{Symmetry:} the reflection $R : x \mapsto 1 - x$
        satisfies $R \partial_{xx} R^{-1} = \partial_{xx}$ but
        $R \partial_x R^{-1} = -\partial_x$. Therefore
        \[
          R L R^{-1} = \kappa \partial_{xx} - \varepsilon \partial_x
          \neq L
        \]
        unless $\varepsilon = 0$.
\end{enumerate}

\subsection{Diagnosis}

The reflection symmetry of the problem and the isotropy requirement
of the diffusion model are incompatible with $\varepsilon \neq 0$.
The violation is detected at the algebraic layer, in the symmetry
criterion, not at the energy layer.

\begin{theorem}[Forced correction]\label{thm:forced-correction}
Under the reflection symmetry constraint, the only model compatible
with the candidate's diffusion hypothesis is $\varepsilon = 0$,
i.e.\ $P_{\rm corr} = P_{\rm cand}|_{\varepsilon = 0}$.
\end{theorem}

\begin{proof}
The candidate $P_{\rm cand}$ is required to be $R$-invariant by the
isotropy hypothesis. But $R L R^{-1} = L$ forces $\varepsilon = 0$.
\end{proof}

\subsection{Discrepancy classification}

The discrepancy is a \emph{structural inconsistency} (regime iii):
$\Delta P$ violates the symmetry required by the problem. The
correction is a \emph{deletion} of the offending term.

\subsection{Remark: non-detection by energy}

The failure of the energy test to detect the error is important. It
illustrates that a \emph{single} layer does not suffice; the
correction arises from the \emph{intersection} of independent
constraints. This is precisely the structural constraining principle
of \S\ref{sec:constraining}.

% ------------------------------------------------------------
\section{Toy 4: Nonlinear $p$-Laplacian with Spectral Emergence}
\label{sec:toy4}
% ------------------------------------------------------------

\subsection{Candidate model}

\[
  P_{\rm cand}: \qquad
  u_t = \partial_x \bigl[ (1 + u_x^2) u_x \bigr],
  \quad x \in (0,1), \quad u(0,t) = u(1,t) = 0.
\]
Equivalently, $u_t = \Ncal(u)$ with
$\Ncal(u) = \partial_x[(1 + u_x^2) u_x]$.

\subsection{Spectral layer: not available directly}

Since $\Ncal$ is nonlinear, no global linear spectral theory applies.
We record
\[
  \Ocal_\Sp = \varnothing
  \qquad \text{(provisionally)}.
\]

\subsection{Geometric layer}

$M = [0,1]$, $g = dx^2$; the operator $\Ncal$ is a nonlinear
generalisation of $-\Delta_g$.

\subsection{Topological layer}

$H_0 = \Z$, $H_1 = 0$; boundary $\partial M = \{0, 1\}$.

\subsection{Algebraic layer: gradient-flow structure}

Define the functional
\[
  E[u] = \int_0^1 \left( \frac{1}{2} u_x^2 + \frac{1}{4} u_x^4 \right) dx.
\]
Its first variation is
\[
  \frac{\delta E}{\delta u}
  =
  -\partial_x(u_x + u_x^3)
  =
  -\Ncal(u).
\]
Therefore, the candidate is the gradient flow
\[
  u_t = -\frac{\delta E}{\delta u},
\]
and consequently
\[
  \frac{dE}{dt} = -\Norm{\Ncal(u)}_{L^2}^2 \leq 0.
\]

\subsection{Numerical layer}

Discretise $E_N(\mathbf u) = h \sum_j [\frac{1}{2}(D_h u_j)^2 +
\frac{1}{4}(D_h u_j)^4]$. The discrete gradient flow is
$\dot{\mathbf u} = -\nabla E_N(\mathbf u)$, preserving the energy
law exactly.

Numerical experiments with $N = 32, 64, 128, \ldots$ show that for
small initial data $\Norm{u_0}_{H^1} = \epsilon \ll 1$, the
solution converges to that of a linear diffusion equation with
coefficients determined by the linearised operator.

\subsection{Spectral emergence}

Compute the Fr\'echet derivative of $\Ncal$ at the equilibrium
$u_\ast = 0$:
\[
  D\Ncal[0] v
  =
  \partial_x(v_x + 3 \cdot 0 \cdot v_x)
  =
  v_{xx}.
\]
Therefore,
\[
  L_0 := D\Ncal[0] = \partial_{xx},
\]
and its spectral data (with Dirichlet boundary conditions) are
\[
  \phi_n(x) = \sin(n \pi x), \qquad \lambda_n = -n^2 \pi^2.
\]

Now the spectral layer \emph{emerges}:
\[
  \Ocal_\Sp^{\mathrm{emg}}
  =
  \bigl( \partial_{xx}, \{-n^2 \pi^2\}_{n \geq 1},
    \{\sin(n \pi x)\}_{n \geq 1} \bigr).
\]

\subsection{Higher-order recovery}

The emergent spectrum captures only the linear part. The higher-order
terms are recovered from the algebraic layer: the functional
$E[u]$ contains $\int u_x^4 dx$, which is not determined by $L_0$
alone.

\begin{proposition}[Non-recovery from spectrum]\label{prop:non-recovery}
The nonlinear term $\int u_x^4 dx$ in $E[u]$ cannot be recovered from
the spectral data of $L_0$.
\end{proposition}

\begin{proof}
Two distinct functionals
\[
  E_1[u] = \int \bigl( \tfrac{1}{2} u_x^2 + \tfrac{1}{4} u_x^4 \bigr) dx,
  \qquad
  E_2[u] = \int \bigl( \tfrac{1}{2} u_x^2 + \tfrac{1}{6} u_x^6 \bigr) dx
\]
have the same second variation at $u = 0$, hence the same Fr\'echet
derivative $L_0 = \partial_{xx}$ and the same spectrum, but different
nonlinear structure.
\end{proof}

\subsection{Reconstruction}

From the emergent spectrum, we reconstruct $L_0 = \partial_{xx}$.
From the algebraic layer, we recover the full functional $E[u]$.
Combining, we reconstruct
\[
  u_t = -\frac{\delta E}{\delta u}
  =
  \partial_x[(1 + u_x^2) u_x]
  =
  P_{\rm cand}.
\]

\subsection{Discrepancy classification}

$\Delta P = 0$ at the reconstruction step: the method closes. But
the example illustrates a new phenomenon: the spectral layer was
\emph{produced} by the method rather than \emph{assumed} at the
start.

\subsection{Validation matrix}

\begin{center}
\begin{tabular}{l|cccccccc}
  \toprule
  Layer & $\mathrm{Ex}$ & $\mathrm{Co}$ & $\mathrm{In}$
  & $\mathrm{Nu}$ & $\mathrm{St}$ & $\mathrm{Cv}$
  & $\mathrm{Re}$ & $\mathrm{Rev}$ \\
  \midrule
  $\Sp$ (emergent) & $\Pmark$ & $\Pmark$ & $\Pmark$ & $\Nmark$ & $\Pmark$ & $\Pmark$ & $\Pmark$ & $\Pmark$ \\
  $\Ge$   & $\Pmark$ & $\Pmark$ & $\Pmark$ & $\Nmark$ & $\Pmark$ & $\Pmark$ & $\Pmark$ & $\Pmark$ \\
  $\To$   & $\Pmark$ & $\Pmark$ & $\Pmark$ & $\Omark$ & $\Pmark$ & $\Pmark$ & $\Pmark$ & $\Pmark$ \\
  $\Al$   & $\Pmark$ & $\Pmark$ & $\Pmark$ & $\Omark$ & $\Pmark$ & $\Pmark$ & $\Pmark$ & $\Pmark$ \\
  $\Nu$   & $\Pmark$ & $\Pmark$ & $\Pmark$ & $\Nmark$ & $\Pmark$ & $\Pmark$ & $\Omark$ & $\Omark$ \\
  $\Fo$   & $\Pmark$ & $\Pmark$ & $\Pmark$ & $\Omark$ & $\Pmark$ & $\Pmark$ & $\Pmark$ & $\Pmark$ \\
  \bottomrule
\end{tabular}
\end{center}

% ============================================================
\part{Validation and Discussion}
% ============================================================

\section{The Validation Matrix for the Toys}
\label{sec:validation-new}

\subsection{Unified Validation Table}

\begin{center}
\renewcommand{\arraystretch}{1.15}
\begin{tabular}{l|c|c|c|c}
  \toprule
  Layer & Toy 1 & Toy 2 & Toy 3 & Toy 4 \\
  \midrule
  $\Sp$   & $\Pmark$ & $\Pmark$ & $\Pmark$ & $\Pmark$ (emergent) \\
  $\Ge$   & $\Pmark$ & $\Pmark$ & $\Pmark$ & $\Pmark$ \\
  $\To$   & $\Pmark$ & $\Pmark$ (bifurcates) & $\Pmark$ & $\Pmark$ \\
  $\Al$   & $\Pmark$ & $\Pmark$ & $\Pmark$ (forces $\varepsilon = 0$) & $\Pmark$ \\
  $\Nu$   & $\Nmark$ & $\Nmark$ & $\Nmark$ & $\Nmark$ \\
  $\Fo$   & $\Pmark$ & $\Pmark$ & $\Pmark$ & $\Pmark$ \\
  $\Cl$   & $\Pmark$ & $\Pmark$ & $\Pmark$ & $\Pmark$ \\
  $\Corr$ & trivial  & trivial  & $\Pmark$ (correction) & trivial \\
  $\An$   & $\Pmark$ & $\Pmark$ & $\Pmark$ & $\Pmark$ \\
  \bottomrule
\end{tabular}
\end{center}

\subsection{Observations}

\begin{enumerate}[label=(\roman*)]
  \item In all four toys, the numerical layer is marked $\Nmark$
        rather than $\Pmark$: numerical evidence is not proof.
  \item Toy~2 exhibits a \emph{bifurcation} at the geometric layer:
        a single discrete object produces two distinct continuous
        models.
  \item Toy~3 exhibits a \emph{correction}: the algebraic layer
        forces $\varepsilon = 0$.
  \item Toy~4 exhibits \emph{spectral emergence}: the spectral layer
        is constructed, not assumed.
\end{enumerate}

\subsection{Constancy of the Validation Structure}

\begin{proposition}[Structural invariance of $\Vmat$]\label{prop:structural-invariance}
Across all four toys, the shape of the validation matrix is
invariant: the numerical and reconstruction rows are always marked
$\Nmark$ and $\Omark$ respectively, while the construction layers
are marked $\Pmark$. This reflects the epistemic asymmetry between
construction and reconstruction.
\end{proposition}

\begin{proof}
The construction layers admit explicit mathematical proofs
(spectral theorem, exterior algebra, Sobolev theory). The numerical
and reconstruction layers depend on the specific discretisation and
inversion schemes, and their proofs are conditional on additional
hypotheses (stability, compactness, regularity).
\end{proof}

% ------------------------------------------------------------
\section{Discussion}
\label{sec:discussion}
% ------------------------------------------------------------

\subsection{Comparison with Related Work}

\paragraph{Structure-preserving numerics.}
The methodology shares with structure-preserving discretisation
\cite{hairer-structure} the goal of preserving qualitative
features of the continuous model at the discrete level. The
distinction is that structure-preserving numerics typically fixes
the continuous model and asks how to discretise it; the present
methodology asks whether the continuous model itself is consistent
with the layer constraints.

\paragraph{Inverse problems.}
The reconstruction pipeline of the framework paper shares the
well-posedness concerns of classical inverse problems
\cite{engl}. The distinction is that the methodology treats
inversion as one of several layers, with the constraining principle
providing a priori restrictions that classical inverse problems
often lack.

\paragraph{Multi-scale and multi-physics modelling.}
The methodology shares the goal of coupling multiple descriptions
\cite{e-weinan}. The distinction is that the coupling here is
\emph{constraining} rather than \emph{informational}: each layer
reduces the admissible set rather than transmitting information to
the next.

\paragraph{Formal mathematics.}
The framework paper's treatment of proof assistants
\cite{mathlib} is inherited here. The toys are simple enough that
their formalisation in Lean~4 would be a reasonable exercise, though
we have not undertaken it in this paper.

\subsection{Limits of the Methodology}

\begin{proposition}[Limits]\label{prop:limits}
The methodology has the following intrinsic limitations:
\begin{enumerate}[label=(\alph*)]
  \item The choice of the initial model class $\mathcal P$ is not
        determined by the method.
  \item The construction of the geometric layer from the spectral
        layer (\S\ref{sec:emergence}) is not canonical: different
        realisations $M$ produce different geometric data.
  \item The emergent spectrum captures only the linear part of a
        nonlinear operator.
  \item The correction procedure (\S\ref{sec:discrepancy}) requires
        an external criterion (symmetry, energy, topology) to
        identify the offending term.
  \item The methodology does not, by itself, produce new
        mathematical theorems; it organises the search for them.
\end{enumerate}
\end{proposition}

\begin{proof}
Each item is an instance of a non-canonicity or information loss
explicit in the corresponding layer definition. The details are
carried by the framework paper's treatment of information loss
(Definition~1.9 of \cite{framework}).
\end{proof}

\subsection{Open Problems}

\begin{enumerate}[label=(\arabic*)]
  \item \textbf{Canonicity of the geometric layer.} Under what
        conditions is there a unique geometric realisation of a
        given spectral datum?
  \item \textbf{Full nonlinear reconstruction.} Can the nonlinear
        part of a gradient flow be recovered from the combination
        of the spectral layer and a finite set of numerical
        moments?
  \item \textbf{Structural constraining dimension.} Given a model
        class $\mathcal P$ and a layer set $\Lam$, does the
        intersection $\Ccal_\ast$ admit a canonical dimension
        (moduli count)?
  \item \textbf{Mechanical verification.} Can the four toys be
        formalised in Lean~4 within Mathlib, and what is the
        shortest proof?
  \item \textbf{Constraining-completeness.} Under what conditions
        does the constraining chain terminate in a singleton?
\end{enumerate}

\subsection{Closing Remark}

The methodology presented in this paper is best described as a
\emph{structural constraining procedure}: it does not produce new
mathematics from nothing, but it organises the interaction between
spectral, geometric, topological, algebraic, numerical, and formal
descriptions of a model in such a way that each description
eliminates a class of candidates and the surviving model is the
intersection of the constraints. The three structural principles
introduced here---constraining, spectral emergence, and discrepancy
classification---give the method its operational content. The four
worked examples show that the method is not vacuous: it detects
symmetry violations, produces corrections, and constructs spectral
data from nonlinear operators where linear theory does not apply.

The central insight, in one line: \emph{in the analysis of a
mathematical model, the question is not what each representation
\emph{shows}, but what each representation \emph{rules out}.}

% ============================================================
\appendix
% ============================================================

\section{Notation Summary}
\label{app:notation}

Inherited from the framework paper \cite{framework}:
\begin{center}
\begin{tabular}{c l}
  \toprule
  Symbol & Meaning \\
  \midrule
  $\Ob$ & Master object \\
  $\Lam$ & Extended layer set \\
  $\Ocal_\lambda$ & Layer-$\lambda$ datum \\
  $\philam$ & Transition morphism \\
  $\rho_\Cl, \rho_\Inv, \rho_\An$ & Reconstruction operators \\
  $\Vmat$ & Validation matrix \\
  $\Pmark, \Nmark, \Omark, \Fmark$ & Proved, numerical, open, failed \\
  \bottomrule
\end{tabular}
\end{center}

New in this paper:
\begin{center}
\begin{tabular}{c l}
  \toprule
  Symbol & Meaning \\
  \midrule
  $\Ccal_\lambda$ & Constrained model space at layer $\lambda$ \\
  $\Ccal_\ast$ & Constrained solution set \\
  $\Ncal$ & Nonlinear operator of evolution equation \\
  $D\Ncal[u_\ast]$ & Fr\'echet derivative at equilibrium $u_\ast$ \\
  $\Ocal_\Sp^{\mathrm{emg}}$ & Emergent spectral layer \\
  $\Delta P$ & Reconstruction discrepancy \\
  \bottomrule
\end{tabular}
\end{center}

\section{Bibliography}
\label{app:biblio}

\begin{thebibliography}{99}

\bibitem{framework}
[Companion framework paper],
\emph{A Multirepresentational Construction--Reconstruction
Methodology for Mathematical Structures}, Parts I--V.

\bibitem{reed-simon-1}
M.~Reed, B.~Simon,
\emph{Methods of Modern Mathematical Physics, Vol.~I: Functional
Analysis}, Academic Press, 1980.

\bibitem{engl}
H.~W.~Engl, M.~Hanke, A.~Neubauer,
\emph{Regularization of Inverse Problems}, Kluwer, 1996.

\bibitem{hairer-structure}
E.~Hairer, C.~Lubich, G.~Wanner,
\emph{Geometric Numerical Integration: Structure-Preserving
Algorithms for Ordinary Differential Equations}, Springer, 2006.

\bibitem{e-weinan}
W.~E, B.~Engquist,
The heterogeneous multiscale methods,
\emph{Comm.~Math.~Sci.}~1 (2003), 87--132.

\bibitem{mathlib}
The Mathlib Community,
\emph{Mathlib: A unified library of mathematics formalized in Lean},
2024.

\bibitem{gordon}
C.~Gordon, D.~Webb, S.~Wolpert,
Isospectral plane domains and surfaces via Riemannian orbifolds,
\emph{Invent.~Math.}~110 (1992), 1--22.

\bibitem{spielman}
D.~Spielman,
\emph{Spectral and Algebraic Graph Theory}, Yale, 2019.

\bibitem{udell}
M.~Udell,
\emph{Introduction to Spectral Graph Theory}, Stanford, 2011.

\end{thebibliography}

\end{document}